\begin{frame}{Antecedentes}
\begin{itemize}
    %    \item La pandemia provocó una caída drástica en el PIB en 2020, pero para 2021 la economía ya se había recuperado.
    \item  De 2019 a 2024 el gasto primario aumentó 3 puntos del PIB, mientras los ingresos apenas si han aumentado.
    %\item El subsidio a los combustibles continúa afectando el desempeño de las finanzas públicas.
    \item A pesar de múltiples reformas tributarias, los ingresos tributarios no parecen ser capaces de superar el 14,5\% del PIB, salvo el repunte transitorio de 2023.
%\end{itemize}
%\end{frame}

%\begin{frame}{Antecedentes}
%\begin{itemize}
    \item A pesar de la activación de la cláusula de escape, los ingresos siguen sobreestimándose.
    \item El deterioro del balance y el crecimiento de la deuda serán más acelerados de lo oficialmente pronosticado.
    \item A medida que transferencias como el FOME y el FEPC se han venido marchitando, otros como el de salud y pensiones han venido aumentando.
\end{itemize}
\end{frame}
